\documentclass[10pt,a4paper]{amsart}
\usepackage{amsmath,amssymb,amsthm,graphicx}
\usepackage{verbatim}
\newenvironment{reference}{\comment}{\endcomment}
%\newenvironment{reference}{}{}
\newenvironment{slogan}{\comment}{\endcomment}
\newenvironment{history}{\comment}{\endcomment}

% For commutative diagrams we use Xy-pic
\usepackage[all]{xy}

% We use 2cell for 2-commutative diagrams.
\xyoption{2cell}
% Theorem environments.
%
\theoremstyle{plain}
\newtheorem{theorem}[subsection]{Theorem}
\newtheorem{proposition}[subsection]{Proposition}
\newtheorem{lemma}[subsection]{Lemma}

\theoremstyle{definition}
\newtheorem{definition}[subsection]{Definition}
\newtheorem{example}[subsection]{Example}
\newtheorem{exercise}[subsection]{Exercise}
\newtheorem{situation}[subsection]{Situation}

\theoremstyle{remark}
\newtheorem{remark}[subsection]{Remark}
\newtheorem{remarks}[subsection]{Remarks}

\numberwithin{equation}{subsection}

% Macros
%
\def\lim{\mathop{\mathrm{lim}}\nolimits}
\def\colim{\mathop{\mathrm{colim}}\nolimits}
\def\Spec{\mathop{\mathrm{Spec}}}
\def\Hom{\mathop{\mathrm{Hom}}\nolimits}
\def\Ext{\mathop{\mathrm{Ext}}\nolimits}
\def\SheafHom{\mathop{\mathcal{H}\!\mathit{om}}\nolimits}
\def\SheafExt{\mathop{\mathcal{E}\!\mathit{xt}}\nolimits}
\def\Sch{\mathit{Sch}}
\def\Mor{\mathop{\mathrm{Mor}}\nolimits}
\def\Ob{\mathop{\mathrm{Ob}}\nolimits}
\def\Sh{\mathop{\mathit{Sh}}\nolimits}
\def\NL{\mathop{N\!L}\nolimits}
\def\CH{\mathop{\mathrm{CH}}\nolimits}
\def\proetale{{pro\text{-}\acute{e}tale}}
\def\etale{{\acute{e}tale}}
\def\QCoh{\mathit{QCoh}}
\def\Ker{\mathop{\mathrm{Ker}}}
\def\Im{\mathop{\mathrm{Im}}}
\def\Coker{\mathop{\mathrm{Coker}}}
\def\Coim{\mathop{\mathrm{Coim}}}

% Boxtimes
%
\DeclareMathSymbol{\boxtimes}{\mathbin}{AMSa}{"02}

%
% Macros for moduli stacks/spaces
%
\def\QCohstack{\mathcal{QC}\!\mathit{oh}}
\def\Cohstack{\mathcal{C}\!\mathit{oh}}
\def\Spacesstack{\mathcal{S}\!\mathit{paces}}
\def\Quotfunctor{\mathrm{Quot}}
\def\Hilbfunctor{\mathrm{Hilb}}
\def\Curvesstack{\mathcal{C}\!\mathit{urves}}
\def\Polarizedstack{\mathcal{P}\!\mathit{olarized}}
\def\Complexesstack{\mathcal{C}\!\mathit{omplexes}}
% \Pic is the operator that assigns to X its picard group, usage \Pic(X)
% \Picardstack_{X/B} denotes the Picard stack of X over B
% \Picardfunctor_{X/B} denotes the Picard functor of X over B
\def\Pic{\mathop{\mathrm{Pic}}\nolimits}
\def\Picardstack{\mathcal{P}\!\mathit{ic}}
\def\Picardfunctor{\mathrm{Pic}}
\def\Deformationcategory{\mathcal{D}\!\mathit{ef}}
\usepackage{cleveref}
\crefname{originallemma}{original lemma}{original lemmas}
\Crefname{originallemma}{Original lemma}{Original lemmas}
\makeatletter
\newlabel{chow-lemma-degrees-and-numerical-intersections}{{0BFI}{source}{}{sourcepointer-1}{}}
\newlabel{chow-lemma-degrees-and-numerical-intersections@cref}{{[originallemma][0][]0BFI}{[1][0][]source}{}{}{}}
\makeatother
\crefname{originalreference}{original reference}{original references}
\Crefname{originalreference}{Original reference}{Original references}
\makeatletter
\newlabel{lemma-pushforward-cap-c1}{{02SU}{source}{}{sourcepointer-2}{}}
\newlabel{lemma-pushforward-cap-c1@cref}{{[originalreference][0][]02SU}{[1][0][]source}{}{}{}}
\makeatother
\crefname{originallemma}{original lemma}{original lemmas}
\Crefname{originallemma}{Original lemma}{Original lemmas}
\makeatletter
\newlabel{spaces-more-morphisms-lemma-chow-noetherian-separated}{{088U}{source}{}{sourcepointer-3}{}}
\newlabel{spaces-more-morphisms-lemma-chow-noetherian-separated@cref}{{[originallemma][0][]088U}{[1][0][]source}{}{}{}}
\makeatother
\crefname{originaldefinition}{original definition}{original definitions}
\Crefname{originaldefinition}{Original definition}{Original definitions}
\makeatletter
\newlabel{spaces-over-fields-definition-intersection-number}{{0EDF}{source}{}{sourcepointer-4}{}}
\newlabel{spaces-over-fields-definition-intersection-number@cref}{{[originaldefinition][0][]0EDF}{[1][0][]source}{}{}{}}
\makeatother
\crefname{originallemma}{original lemma}{original lemmas}
\Crefname{originallemma}{Original lemma}{Original lemmas}
\makeatletter
\newlabel{spaces-over-fields-lemma-intersection-number-and-pullback}{{0EDJ}{source}{}{sourcepointer-5}{}}
\newlabel{spaces-over-fields-lemma-intersection-number-and-pullback@cref}{{[originallemma][0][]0EDJ}{[1][0][]source}{}{}{}}
\makeatother
\crefname{originallemma}{original lemma}{original lemmas}
\Crefname{originallemma}{Original lemma}{Original lemmas}
\makeatletter
\newlabel{spaces-over-fields-lemma-numerical-polynomial-leading-term}{{0EDE}{source}{}{sourcepointer-6}{}}
\newlabel{spaces-over-fields-lemma-numerical-polynomial-leading-term@cref}{{[originallemma][0][]0EDE}{[1][0][]source}{}{}{}}
\makeatother
\title{Intersection numbers agree with the degree of the iterated Chern class cap}
\author{The Stacks Project Authors}
\date{stacks-project:spaces-chow.tex:local-snapshot-20261012; GNU FDL 1.2 or later}
\begin{document}
\maketitle
\noindent\textit{Editorial note.} All delivered mathematical text is the author's own, copied from the pinned Stacks Project chapter file spaces-chow.tex; no statement or proof step was rewritten, reordered or removed.\par\medskip
\noindent\textit{Editorial note.} The counted proof is the author's complete argument: a reduction to an integral space by multiplicities, a passage to a proper scheme through Chow's lemma, and a check that both sides survive the replacement\par\medskip
\noindent\textit{Editorial note.} The label lemma-degrees-and-numerical-intersections names the counted statement, and the proof environment immediately following it in the source is delivered verbatim as its proof.\par\medskip
\noindent\textit{Editorial note.} The counted statement is carried by the author's own lemma environment in the chapter source.\par\medskip
\noindent\textit{Allowed standard prerequisites (editorial).}\begin{itemize}
\item The definition of the intersection number of invertible sheaves with an integral closed subspace, and its agreement with the leading term of the corresponding numerical polynomial
\item The decomposition of a cycle into the classes of its irreducible components with multiplicities
\item Chow's lemma producing a proper morphism that is an isomorphism over a dense open subspace and whose source is a scheme, the compatibility of intersection numbers with pullback, and the projection formula for capping with a first Chern class
\end{itemize}\medskip
\ifdefined\hypertarget\hypertarget{sourcepointer-1}{}\fi
\noindent\textit{Source reference.} Original lemma 0BFI: sources/stacks-chow/tags.txt, source lines 14386--14386. Prerequisite pointer invoked at proof line(s) 5264; it names a standard fact of a sibling chapter that the counted argument uses, and every step of the counted proof built on it is retained verbatim in the delivered body.\par\medskip
\ifdefined\hypertarget\hypertarget{sourcepointer-2}{}\fi
\noindent\textit{Source reference.} Original reference 02SU: sources/stacks-spaces-chow/pinned-original-source.tex, source lines 2933--2933. Prerequisite pointer invoked at proof line(s) 5261; it names a standard fact of the same chapter that the counted argument uses, and every step of the counted proof built on it is retained verbatim in the delivered body.\par\medskip
\ifdefined\hypertarget\hypertarget{sourcepointer-3}{}\fi
\noindent\textit{Source reference.} Original lemma 088U: sources/stacks-chow/tags.txt, source lines 10474--10474. Prerequisite pointer invoked at proof line(s) 5246; it names a standard fact of a sibling chapter that the counted argument uses, and every step of the counted proof built on it is retained verbatim in the delivered body.\par\medskip
\ifdefined\hypertarget\hypertarget{sourcepointer-4}{}\fi
\noindent\textit{Source reference.} Original definition 0EDF: sources/stacks-chow/tags.txt, source lines 17990--17990. Prerequisite pointer invoked at proof line(s) 5225; it names a standard fact of a sibling chapter that the counted argument uses, and every step of the counted proof built on it is retained verbatim in the delivered body.\par\medskip
\ifdefined\hypertarget\hypertarget{sourcepointer-5}{}\fi
\noindent\textit{Source reference.} Original lemma 0EDJ: sources/stacks-chow/tags.txt, source lines 17994--17994. Prerequisite pointer invoked at proof line(s) 5255; it names a standard fact of a sibling chapter that the counted argument uses, and every step of the counted proof built on it is retained verbatim in the delivered body.\par\medskip
\ifdefined\hypertarget\hypertarget{sourcepointer-6}{}\fi
\noindent\textit{Source reference.} Original lemma 0EDE: sources/stacks-chow/tags.txt, source lines 17989--17989. Prerequisite pointer invoked at proof line(s) 5238; it names a standard fact of a sibling chapter that the counted argument uses, and every step of the counted proof built on it is retained verbatim in the delivered body.\par\medskip
\begin{lemma}
\label{lemma-degrees-and-numerical-intersections}
Let $k$ be a field. Let $X$ be a proper algebraic space over $k$.
Let $Z \subset X$ be a closed subspace of dimension $d$.
Let $\mathcal{L}_1, \ldots, \mathcal{L}_d$ be invertible
$\mathcal{O}_X$-modules. Then
$$
(\mathcal{L}_1 \cdots \mathcal{L}_d \cdot Z) =
\deg(
c_1(\mathcal{L}_1) \cap \ldots \cap c_1(\mathcal{L}_1) \cap [Z]_d)
$$
where the left hand side is defined in
Spaces over Fields, Definition
\ref{spaces-over-fields-definition-intersection-number}.
\end{lemma}

\begin{proof}
Let $Z_i \subset Z$, $i = 1, \ldots, t$ be the irreducible components
of dimension $d$. Let $m_i$ be the multiplicity of $Z_i$ in $Z$. Then
$[Z]_d = \sum m_i[Z_i]$ and
$c_1(\mathcal{L}_1) \cap \ldots \cap c_1(\mathcal{L}_d) \cap [Z]_d$
is the sum of the cycles
$m_i c_1(\mathcal{L}_1) \cap \ldots \cap c_1(\mathcal{L}_d) \cap [Z_i]$.
Since we have a similar decomposition for
$(\mathcal{L}_1 \cdots \mathcal{L}_d \cdot Z)$ by
Spaces over Fields, Lemma
\ref{spaces-over-fields-lemma-numerical-polynomial-leading-term}
it suffices to prove the lemma in case $Z = X$
is a proper integral algebraic space over $k$.

\medskip\noindent
By Chow's lemma there exists a proper morphism $f : X' \to X$
which is an isomorphism over a dense open $U \subset X$
such that $X'$ is a scheme. See More on Morphisms of Spaces, Lemma
\ref{spaces-more-morphisms-lemma-chow-noetherian-separated}.
Then $X'$ is a proper scheme over $k$. After replacing $X'$
by the scheme theoretic closure of $f^{-1}(U)$
we may assume that $X'$ is integral. Then
$$
(f^*\mathcal{L}_1 \cdots f^*\mathcal{L}_d \cdot X') =
(\mathcal{L}_1 \cdots \mathcal{L}_d \cdot X)
$$
by Spaces over Fields, Lemma
\ref{spaces-over-fields-lemma-intersection-number-and-pullback}
and we have
$$
f_*(c_1(f^*\mathcal{L}_1) \cap \ldots \cap c_1(f^*\mathcal{L}_d) \cap [Y]) =
c_1(\mathcal{L}_1) \cap \ldots \cap c_1(\mathcal{L}_d) \cap [X]
$$
by Lemma \ref{lemma-pushforward-cap-c1}. Thus we may replace $X$ by $X'$
and assume that $X$ is a proper scheme over $k$. This case
was proven in Chow Homology, Lemma
\ref{chow-lemma-degrees-and-numerical-intersections}.
\end{proof}

\end{document}
